\begin{proof}
La afirmación (i) ya se vio durante la demostración de II) y no
depende de la hipótesis sobre la dimensión de $A$. La hipótesis de que
el morfismo $\operatorname{Spec}(A')\to\operatorname{Spec}(A)$ es normal
implica que $A'\otimes_A k=A'/\mathfrak mA'$ es un anillo normal \emph{local} y
noetheriano, por tanto \emph{íntegro}; esto muestra ya
que $\mathfrak p'=\mathfrak mA'$ es primo. Según (6.1.2), se tiene
$\dim(A'_{\mathfrak p'})=\dim(A)+
\dim(A'_{\mathfrak p'}/\mathfrak mA'_{\mathfrak p'})$. Pero
$A'_{\mathfrak p'}/\mathfrak mA'_{\mathfrak p'}=
A'_{\mathfrak p'}/\mathfrak p'A'_{\mathfrak p'}$ es el cuerpo
residual de $A'_{\mathfrak p'}$, por lo que
$\dim(A'_{\mathfrak p'})=1$. El hecho de que $A'_{\mathfrak p'}$
esté contenido en el producto directo de los $L'_j$ tales que
$\mathfrak q'_j\subset\mathfrak p'$ se debe a que los ideales
$\mathfrak q'_jA'_{\mathfrak p'}$ para esos índices $j$ son los
ideales primos mínimos de $A'_{\mathfrak p'}$.
